%%
%% Ergänzung zu Kapitel 6: iOS-/iPadOS-Studienumsetzung
%%

\begin{neu}
\section{Plattformübertragung der Studienintervention auf iOS und iPadOS}

Die native Apple-Portierung wurde mit dem Ziel umgesetzt, die experimentell relevanten Abläufe der Android-Anwendung ohne Einführung einer zusätzlichen Studienbedingung abzubilden. Der innerhalb der App realisierte Ablauf bleibt deshalb unverändert als Within-Subject-Sequenz mit zehn Cue-Matching- und anschließend zehn Cue-Labeling-Situationen definiert. Jede Situation enthält fünf Trials und wird durch genau einen Craving-Selbstbericht abgeschlossen. Bei vollständiger Durchführung entstehen somit weiterhin 100 Trials und 20 wissenschaftliche Selbstberichte pro teilnehmender Person.

Für die Plattformübertragung wurden die kanonischen Reize und Labelzuordnungen übernommen. Trialauswahlen und Reaktionszeiten werden auch unter iOS nicht Teil des wissenschaftlichen Datensatzes. Der Submission-Request enthält keine Plattformkennung und keine vom Client bestimmte Versuchsbedingung; die Bedingungszuordnung erfolgt weiterhin serverseitig anhand der Zahl bereits bestätigter Berichte. Damit entsteht auf Ebene des Forschungsdatensatzes dieselbe Struktur wie bei Android: eine pseudonyme Teilnehmerkennung, der serverseitig bestimmte Bedingungscode und der Craving-Wert.

Plattformspezifisch ist lediglich die konkrete native Ausführung einzelner UI- und Zufallsmechanismen. Die 50 Matching-Items werden auf der Apple-Plattform einmalig zufällig permutiert und in Blöcke zu fünf Items aufgeteilt; auch die Position der Auswahlmöglichkeiten wird je Trial zufällig festgelegt. Diese konkrete Zufallsfolge kann sich zwischen Android und iOS unterscheiden, verändert aber weder Zahl noch Inhalt der vorgesehenen Stimuli oder die Reihenfolge der beiden Interventionsphasen. Eine solche technische Zufallsrealisierung ist daher nicht als zusätzliche experimentelle Bedingung zu interpretieren.

Auch die zeitliche Freigabelogik bleibt erhalten. Nach einem erfolgreich bestätigten Selbstbericht wird im Release-Build eine dreistündige Sperrzeit bis zur nächsten Situation gesetzt. Ein noch ausstehender Submission- oder Abschlusszustand blockiert einen neuen Durchlauf vollständig. Der lokale Fortschritt wird erst nach einer gültigen Serverantwort erhöht. Dadurch soll verhindert werden, dass eine nur lokal als erfolgreich angenommene Übertragung zu einer Verschiebung der Reihenfolge oder zu einem fehlenden Selbstbericht führt.

Für die Rekrutierungskanäle gelten dieselben beiden Abschlussmodi wie auf Android. Direkt Teilnehmende erhalten nach dem zwanzigsten Bericht den CueLens-Abrechnungscode mit separater Bestätigung; Prolific-Teilnehmende erhalten keinen solchen Code und werden über den serverseitigen \texttt{PROLIFIC\_MANUAL}-Pfad abgeschlossen. Die Apple-App verändert damit weder das Vollständigkeitskriterium von 20 Selbstberichten noch die administrative Zuordnung der Vergütung.

Methodisch muss zwischen \emph{implementierter Plattformäquivalenz} und \emph{tatsächlicher Nutzung in der Stichprobe} unterschieden werden. Die Implementierung belegt, dass der Studienablauf technisch auch auf iPhone und iPad abbildbar ist. Daraus folgt jedoch nicht, dass bereits erhobene Android-Daten gemeinsam mit iOS-Daten ausgewertet werden dürfen oder dass die Apple-Plattform Bestandteil einer bereits abgeschlossenen Rekrutierung war. Für die Ergebnisdarstellung ist deshalb ausschließlich die Plattformverteilung der tatsächlich eingeschlossenen und ausgewerteten Teilnehmenden maßgeblich. Wurde die Apple-Version erst nach Abschluss oder außerhalb des freigegebenen Erhebungszeitraums fertiggestellt, ist sie als nachträgliche technische Erweiterung des Prototyps zu dokumentieren.

Zum dokumentierten Entwicklungsstand ist die iOS-/iPadOS-Implementierung funktional vollständig, aber noch nicht formal zur Distribution freigegeben. Das Definition-of-Done-Audit weist reale Staging-Ende-zu-Ende-Prüfungen für direkte und Prolific-Registrierungen, TestFlight-Tests auf mehreren realen Geräten, Distribution-Entitlements, den Langzeittest sowie externe Datenschutz- und Ethiknachweise als noch offen aus. Für die Studienmethodik bedeutet dies, dass die technische Äquivalenz durch Implementierung und automatisierte Prüfungen gestützt wird, ein produktiver Einsatz in einer laufenden Studie jedoch zusätzlich die organisatorische und formale Freigabe erfordert.
\end{neu}
